\documentclass[11pt,a4paper]{article}
\usepackage[T1]{fontenc}
\usepackage[utf8]{inputenc}
\usepackage[margin=27mm]{geometry}
\usepackage{amsmath}
\usepackage{booktabs,tabularx}
\usepackage{listings}
\usepackage{xcolor}
\usepackage{xurl}
\usepackage{hyperref}
\hypersetup{colorlinks=true,linkcolor=blue!45!black,urlcolor=blue!45!black,citecolor=blue!45!black}
\urlstyle{tt}
\lstset{basicstyle=\small\ttfamily,columns=fullflexible,keepspaces=true,
breaklines=true,frame=single}
\setlength{\emergencystretch}{3em}
\newcommand{\code}[1]{\nolinkurl{#1}}
\newcommand{\src}[1]{\path{#1}}
\newcommand{\evidence}[1]{\par\smallskip{\small\raggedright\noindent\textit{Implementation evidence:} #1\par}}
\title{Ibex as an Analytics Engine\\\large Data representation, planning, and execution
\\Working technical description}
\author{Bob Jansen, OpenAI GPT}
\date{\today}

\begin{document}
\maketitle
\begin{abstract}
  \noindent
  Ibex exposes analytical operations through a statically typed language and
  executes them over columnar data. This document explains the implementation
  boundary between that language and its host analytics runtime: how table data
  is represented, how logical operations become executable operators, where
  intermediate results are stored, and how one query uses multiple threads.
  It connects these mechanisms to query-processing literature while keeping
  literature comparisons separate from claims about Ibex. This first draft is
  deliberately incomplete: further operator specializations and design rationale
  remain subjects for investigation and author review.
\end{abstract}

\section{Scope and evidence}
The primary audience is external database-systems readers. The text should
explain the design and its relationship to established techniques, with enough
operational detail to follow how a query executes. It assumes familiarity with
basic relational operations and programming, but not the Ibex implementation.
\src{SPEC.md} remains the authority
for language semantics. This document describes mechanisms behind those
semantics; it does not introduce syntax or promise performance.

Three kinds of statement are distinguished throughout. \emph{Implementation
evidence} names code and symbols inspected for a claim. \emph{Literature
context} explains related published techniques; similarity does not establish
historical influence or equivalent implementation. \emph{Open questions} name
unfinished investigations. Existing runtime notes are useful context, but
their architectural aspirations are not treated as proof that every execution
path satisfies them.

\section{The execution architecture}
The host runtime is a single-process, shared-memory engine. Its public
\code{interpret} entry point claims a query-execution lease, constructs an
operator, and drains it into a materialized result table. A concurrent or
re-entrant top-level call is rejected with an error. This is a restriction on
the host runtime entry point; it does not imply that a query executes on one
thread, or specify how an embedding application manages multiple processes.

Interpreted and compiled execution are functionally equivalent implementations
of the same Ibex language. Both must implement the same specified behavior;
a functional divergence is a bug unless explicitly designated as an intentional
exception. As of writing there are no intentional exceptions. Different
planning, scheduling, and code-generation mechanisms do not establish different
language semantics or feature sets. This is the required contract, not a claim
that the implementation is currently bug-free.

There are two execution routes. The interpreted route lowers input into an
intermediate representation (IR) and evaluates it through the runtime. A
separate emitter translates IR into C++ source, including calls to
\code{ibex::ops} functions. The distinction below explains how each route
implements the shared functionality.

\begin{center}
  \begin{tabular}{c}
    Ibex source\\
    $\downarrow$ parsing, analysis, lowering\\
    logical IR and rewrites\\
    $\swarrow$ \hspace{5em} $\searrow$\\
    \begin{tabular}{c@{\hspace{3em}}c}
      host execution & C++ source emission\\
      physical planning / fallback & external C++ compilation\\
      operators and kernels & generated executable\\
      materialized result & calls into C++ operations
    \end{tabular}
  \end{tabular}
\end{center}
This outline permits different sequences of analyses while requiring
equivalent functionality. Rejections encountered in the emitter are recorded
as parity bugs in Appendix~\ref{sec:backend-bugs}. An ``unsupported'' diagnostic
documents an implementation gap; it does not by itself designate an intentional
exception to equivalence.

\evidence{\src{src/runtime/interpreter.cpp}, \code{interpret};
\src{include/ibex/runtime/query_lease.hpp}, \code{QueryExecutionLease};
\src{tools/ibex_compile.cpp}; \src{src/codegen/emitter.cpp},
\code{Emitter::emit_node}, especially Filter, Project, and Window cases.}

\paragraph{Literature context.}
Neumann's query-compilation work describes data-centric code generation that
pushes tuples through generated pipelines and aims to retain values in CPU
registers~\cite{neumann2011}. Emitting C++ does not by itself establish that
Ibex implements that execution model. The supported statement here is that
Ibex has an IR-to-C++ backend alongside its host operator runtime.

\section{Data representation and ownership}
Ibex's column layouts are heavily inspired by the Apache Arrow columnar
format~\cite{arrowColumnar}. The relevant Arrow layouts are fixed-size
primitive arrays, variable-size binary arrays (including UTF-8 strings),
and dictionary-encoded arrays, together with separate validity bitmaps.
The following describes how these ideas appear in Ibex's internal storage.

\subsection{Tables contain typed columns}
A runtime \code{Table} contains an ordered vector of \code{ColumnEntry}
objects and a name-to-position index. Each entry carries a name, shared
ownership of a \code{ColumnValue}, and optional validity information.
\code{ColumnValue} is a variant of concrete typed columns: 64-bit integers,
double-precision floating point, strings, categoricals, dates, timestamps,
and booleans in the inspected header. A statically typed language therefore
coexists with runtime dispatch over concrete column representations.

The useful physical distinction is between an element's logical meaning and
its storage. In particular, strings can use either a flat byte representation
or dictionary encoding.

\begin{center}
  \begin{tabularx}{\linewidth}{@{}lX@{}}
    \toprule
    Representation & Layout in the inspected implementation\\
    \midrule
    Fixed-width & Contiguous typed values; ordinary owned storage uses a vector.\\
    Flat strings & One byte buffer and $n+1$ unsigned 32-bit offsets for $n$ rows.\\
    Categorical strings & Signed 32-bit codes and a shared string dictionary.\\
    Booleans & Values packed into 64-bit words.\\
    Validity & An optional bitmap, separate from the value representation.\\
    \bottomrule
  \end{tabularx}
\end{center}
For a flat string column, row $i$ occupies bytes
$[o_i,o_{i+1})$. For a categorical column, its value is $D[c_i]$, where $D$
is the dictionary. The same integer code in two unrelated dictionaries need
not mean the same string. This matters when combining chunks or comparing
encoded keys; a code is meaningful only together with its dictionary.
These layouts correspond to Arrow's ``Variable-size Binary Layout'' and
``Dictionary-encoded Layout''~\cite{arrowColumnar}. Arrow's conventional
binary/string layout uses signed 32-bit offsets (or signed 64-bit offsets
for its large variant), whereas the Ibex storage described above uses
unsigned 32-bit offsets. The shared design therefore does not imply that
every internal buffer is interchangeable without checking its representation.

\subsection{Nulls and sharing}
An absent validity bitmap means every row is valid. With a bitmap, a set bit
means valid and a cleared bit means null. This follows Arrow's validity
convention, including permission to omit the bitmap when there are no
nulls~\cite{arrowColumnar}. Payload storage and validity are
distinct: inspecting a stored numeric payload alone does not establish
whether that row contains a null. The full semantics of null propagation in
expressions and aggregates remain to be documented separately.

Column entries share their value objects through \code{shared_ptr}. The
\code{Chunk::mutable_column} accessor copies the value object if it is
shared before exposing mutable access. The inspected generic column also
supports immutable external buffers with a shared lifetime owner and detaches
on mutation. These are specific ownership mechanisms, not a blanket guarantee
that arbitrary queries or interchange conversions are zero-copy.

\evidence{\src{include/ibex/runtime/interpreter.hpp}, \code{ColumnValue},
\code{ColumnEntry}, \code{ValidityBitmap}, and \code{Table};
\src{include/ibex/core/column.hpp}, the generic \code{Column} and its string,
categorical, and boolean specializations;
\src{include/ibex/runtime/operator.hpp}, \code{Chunk::mutable_column}.}

\subsection{Row layout is part of the execution contract}
Tables also carry derived properties describing ordering, a time index, and
grouping. \code{TableProperties} centralizes their construction and transfer.
Its \code{RowTransform} classification distinguishes preserving rows,
selecting a subsequence, reordering rows, and recombining rows. Preserving a
subsequence can retain ordering; constructing rows from groups or multiple
inputs cannot simply inherit every input claim.

The time-frame constructor establishes ascending order on its time index.
Normalization and key-fate rules constrain what survives later operations.
Thus, time-series execution depends on facts about row layout as well as on
timestamp values. A detailed account of window boundaries, group boundaries,
and ordering-sensitive functions is reserved for a subsequent chapter.

\evidence{\src{include/ibex/runtime/table_properties.hpp},
\code{RowTransform}, \code{TableProperties::derive},
\code{TableProperties::time_frame}, and \code{normalized}.}

\section{From logical operations to executable work}
\subsection{Lowering and logical rewrites}
The whole-program \code{parser::lower} entry performs effect analysis,
lowers the program, validates column references against known schemas, and
checks joins. It then applies schema-aware join filter pushdown, semi-join
pushdown, and inner-to-semi reductions before invoking the default optimizer.
The default pass manager removes eligible unused pure preamble calls,
canonicalizes the IR, and annotates pending order requirements.

Two limits to this statement matter. First, static checks use the schemas
available at that boundary; this is not evidence that all data-source schemas
are known at parse time. Second, some analyses live in the calling path.
The REPL also computes required columns and scan predicates for lazy-source
decoding. A list of passes from \code{optimizer.cpp} alone would consequently
be an incomplete description of optimization.

Effect summaries constrain which calls may be discarded or reordered.
\code{is_reorderable}, for example, checks non-I/O effects and conflicts
between declared I/O resources. The implementation therefore represents
semantic constraints beyond relational schema when deciding whether a
transformation is permitted.

\evidence{\src{src/parser/lower.cpp}, \code{lower};
\src{src/ir/optimizer.cpp}, \code{make_default_pass_manager},
\code{is_elidable}, and \code{is_reorderable};
\src{src/repl/repl.cpp}, calls to \code{required_columns},
\code{scan_predicates}, and \code{decode_demanded_lazy_sources}.}

\subsection{Physical planning and fallback}
At runtime, \code{build_operator_impl} calls \code{physical::plan_physical}.
The resulting description can cover a chain of row-local operations or a
supported stateful operator. The physical plan records such choices as map
steps, fusion, source kind, parallel mode, and join or aggregate strategy.
For a migrated plan, the executor consumes this description. Other shapes
use explicit handling or a materialized fallback through
\code{build_materialized_fallback}.

Fallback is part of the current architecture. It lets supported pipelines
and operators coexist with operations that consume whole tables. In
particular, falling back for a parent need not discard physical execution of
its children: the fallback builder routes direct children through the
physical path. A precise account must state both the preferred strategy and
the conditions that select another strategy.

For joins, the plan distinguishes streaming probe from materializing both
inputs. Its decline reasons include unsupported key shapes, non-equi
predicates, null-equality semantics, cardinality assertions, and match
selection. Textual left and right inputs are recorded separately from
runtime build/probe orientation. For aggregates, the strategy vocabulary
includes an adaptive sorted-stream path with a hash fallback, a fused
left-join count, and whole-input materialization. These are strategy
boundaries; their detailed algorithms have not yet been audited for this
document.

\evidence{\src{src/runtime/runtime_entry.cpp}, \code{build_operator_impl};
\src{src/runtime/physical_plan.hpp}, \code{Plan}, \code{JoinPlan},
\code{JoinDeclineReason}, and \code{AggregateStrategy};
\src{tests/test_physical_plan.cpp}, tests for map fusion, join classification,
aggregate classification, and pipelines over breakers.}

\section{Chunks, pull execution, and materialization}
The common operator interface is
\begin{lstlisting}
next() -> expected<optional<Chunk>, string>
\end{lstlisting}
A successful pull produces a chunk or signals end of input; an error carries
a message. A chunk is a horizontal portion of a table, represented by its
columns rather than by an array of row objects. It also has sequence and
source-row-offset fields. A chunk is an interface unit, not a guarantee of
a particular size: \code{TableSourceOperator} can emit an entire materialized
table in one chunk.

\code{MaterializeOperator::run} pulls chunks and assembles a \code{Table}.
The first chunk donates its columns; later chunks are appended with schema
checks and validity handling. The runtime contract requires an empty input
to emit an empty chunk carrying its schema, so zero rows do not erase the
result's column names and types. A dedicated physical-execution test checks
this behavior for an empty filtered input.

Incrementally draining the final sink avoids retaining every incoming chunk
as a separate object. It does \emph{not} establish bounded memory for the
whole query: the final result, retained inputs, hash tables, sort state,
buffered pipeline results, and materialized fallback inputs may also occupy
memory. This draft makes no claim about spilling to disk or a global memory
limit.

\evidence{\src{include/ibex/runtime/operator.hpp}, \code{Operator},
\code{Chunk}, \code{TableSourceOperator}, and \code{MaterializeOperator};
\src{src/runtime/CONTRACTS.md}, chunk and materialization contracts;
\src{tests/test_physical_plan.cpp}, ``Migrated filter keeps the empty input's
schema carrier''.}

\paragraph{Literature context.}
Volcano provides the classic reference for composable operators with a
standard iterator interface~\cite{graefe1994}. Ibex's pull interface is
comparable at that architectural level, but its transfer unit is a columnar
chunk. MonetDB/X100 motivates vector-at-a-time processing as a way to
amortize interpretation overhead while avoiding full-column intermediates
at every step~\cite{boncz2005}. Ibex's chunk and kernel boundaries provide a
useful comparison, without establishing identical vector sizes, instruction
generation, or performance. Here, ``vector'' means a batch of values; it
does not by itself assert use of hardware SIMD instructions.

\section{A small query through the engine}
Consider the following query over an already registered table:
\begin{lstlisting}
trades[filter price > 15, select { price }];
\end{lstlisting}
The physical-plan test for this exact expression establishes that
\code{Project(Filter(Scan(trades)))} becomes one fused map step with
\code{FilterProjectGather} capability and a \code{TableScan} source.
Projection is attached to the filter step, so the physical execution need
not preserve two independent operators merely because the logical tree
contains them.

For illustration, suppose the non-null price column is $(10,20,15,30)$.
The predicate's survivor positions are $(1,3)$ under zero-based indexing,
and the selected output column is $(20,30)$. The mathematical effect is
\[
  S = (i \mid 0\leq i<n,\ p_i>15),\qquad
  r_j=p_{S_j}.
\]
These positions explain the gather operation; this notation does not claim
that every execution path allocates a survivor-index vector.

The chosen implementation can execute serially or, when eligible, use
parallel range processing. To share the work, the engine divides the input
into small ranges of consecutive rows called \emph{morsels}. Each morsel is
a portion of work that a worker can process independently; for example,
one worker can filter an earlier range while another filters a later range.
A morsel describes which input rows to work on, while a chunk is the data
passed between operators. Processing a morsel can produce a chunk containing
only its surviving rows.

A separate two-phase filter specialization first
counts survivors in each morsel, computes prefix offsets, and then writes
disjoint output slices. If morsel $k$ has $c_k$ survivors, its output begins
at $b_k=\sum_{j<k}c_j$. This separates completion order from result position:
the second morsel may finish first and still write after the first morsel's
reserved output. This is an explanation of the specialization, not a claim
that the tiny illustrative input passes its parallel thresholds.

\evidence{\src{tests/test_physical_plan.cpp}, ``Physical plan lowers
filter+select into a fused map step'';
\src{src/runtime/pipeline_executor.cpp}, \code{TwoPhaseFilterOperator}
and the specialization's selection in pipeline construction.}

\section{Parallel execution within one query}
Three mechanisms should be kept distinct when reasoning about concurrency.
\emph{Operator-internal fan-out} divides work inside an operation and joins
the tasks before that phase completes. A \emph{morsel pipeline} carries each
small input range through a chain of eligible row-local operations, such as
filtering followed by computing a column. Workers process these ranges of a
materialized input independently, and their outputs are merged in sequence.
\emph{Pipeline overlap} allows scan
or child-stage production to overlap consumption through bounded buffers.

The worker pool uses a shared task queue and completion state. In the morsel
executor, workers obtain range sequence numbers from an atomic cursor.
These are separate levels: the pool runs worker bodies, while the cursor
distributes morsels among those bodies. The inspected morsel pipeline owns
its materialized input before worker execution. Deferred scan pipelines have
a separate source-unit mechanism; this document does not conflate them with
partitioning an already resident table.

\code{PipelinedScanOperator} processes source units and makes results
available in sequence. \code{PipelinedStageOperator} uses a producer thread
and bounded buffering. The latter is distinct from a short compute task
submitted to the worker pool. Parallelism policy also distinguishes a
compute budget from the number of available decode workers; the complete
configuration and admission rules remain to be catalogued.

\evidence{\src{src/runtime/worker_pool.cpp}, the task queue,
\code{WorkerPool::submit}, and \code{wait_for_batch};
\src{src/runtime/pipeline_executor.cpp}, \code{MorselPipelineOperator},
\code{PipelinedScanOperator}, and \code{PipelinedStageOperator};
\src{src/runtime/PARALLELISM.md}, the three layers and thread budgets.}

\paragraph{Literature context.}
Leis et al. describe scheduling small input fragments through operator
pipelines, with NUMA locality and elastic allocation between
queries~\cite{leis2014}. Ibex's morsel pipeline shares the fragment-through-a-
pipeline idea. Its host query lease and shared-cursor scheduling are more
specific constraints. The word ``morsel'' here does not establish NUMA-aware
placement, work stealing, or elastic allocation among concurrent queries.

\subsection{Determinism requires a scoped statement}
The design priority is to preserve the row-order guarantees needed to reason
about time-series computations. Exact floating-point reproducibility is a
separate requirement and may be relaxed; doing so does not relax output row
ordering. The permitted numerical differences and their scope remain to be
specified in a later revision. This direction does not itself change the
current language specification or establish a numerical tolerance.

Ordered merges and prefix-assigned output ranges make row position
independent of task completion order. The runtime contracts set a strong
serial/parallel equivalence goal, but also acknowledge floating-point
reduction exceptions. Consequently this draft does not promote the broad
``byte-identical at any core count'' wording in runtime notes into an
unqualified engine guarantee.

For reductions, stable output order and numerical reproducibility are
different properties: in floating-point arithmetic, $(a+b)+c$ can differ
from $a+(b+c)$. A reproducibility claim must specify the partitioning and
merge order, the comparison target, and the affected aggregates. Error
selection and cancellation also require their own examination. Existing
serial/parallel tests provide evidence for particular inputs and paths,
not a proof covering the complete engine.

\evidence{\src{src/runtime/CONTRACTS.md}, determinism contract and exceptions;
\src{tests/test_physical_plan.cpp}, ``Migrated pipelines produce identical
results in serial and parallel''.}

\section{Operational detail: sources, operators, and retained state}
\label{sec:operations}
The preceding sections describe the engine's components. This section follows
their execution boundaries: when a source is read, what an operator retains
between pulls, and when output becomes available. These concrete mechanisms
describe the host runtime; the compiled route must provide equivalent
functionality under the contract in Section~2.

\subsection{Binding a source and deciding what to read}
A registered materialized table already owns or shares its column buffers.
A lazy table instead starts with a schema, a row count, and a decoder supplied
by its source. This allows a caller to resolve column demand before asking for
column values. \code{LazyTable::project} finds demanded columns missing from
its cache, decodes them, and returns columns in source-schema order. Its cache
contains whole-source columns, so subsequent uses can reuse those buffers.

Filtering adds a second decision: which rows need payload values? In the
ordinary \code{project_where} path, the engine obtains columns needed by
predicates, evaluates the predicates, and constructs a selection of surviving
source-row indices. It then requests the demanded payload columns for those
rows. Already cached columns can be gathered in memory. Newly decoded selected
payloads are not inserted as whole-source cache entries: otherwise a later
query could mistake a filtered subset for the original column.

For example, a scan that filters ticks on \code{price} and returns only
\code{ts} and \code{volume} needs the predicate's answer but does not need
\code{price} in its output. An eligible reader can evaluate certain range or
string predicates while decoding, returning the selection without constructing
the full predicate column. If the reader reports that it cannot answer that
request, the engine uses ordinary decode-and-filter execution. An error is
different from a declined optimization and propagates to the caller.

This separates three quantities: columns referenced by predicates, columns
required by consumers, and source rows selected for those consumers. It also
explains why column-demand analysis and decoding cannot be understood solely
by looking at the final projection.

\evidence{\src{src/runtime/lazy_table.cpp}, \code{LazyTable::project},
\code{project_where}, \code{decode_whole_columns}, and \code{project_rows};
\src{include/ibex/runtime/lazy_table.hpp}, \code{ColumnDecodeFn} and
\code{LazySourceReader}.}

\subsection{Source units and reader ownership}
A lazy source may expose a sequence of \code{SourceUnit} ranges that covers
its rows in order. A reader without such a decomposition is read as a whole.
For a decomposed source, a unit restricts decoding to part of the source;
selection indices still refer to the original source, not to positions
starting at zero within that unit. The reader intersects the selection with
the unit. Keeping one index space prevents a selection made before decoding
from selecting different rows after partitioning.

Each \code{LazySourceReader} owns its mutable decoder state. The reader
factory and pool permit independent readers to share immutable resources
without sharing a decoder cursor. In a scan pipeline, workers can produce
different source units concurrently, while the sequence-indexed output buffer
controls when the consumer receives them. Buffer capacity applies
backpressure: production must wait when the consumer has not released space.
The unit size belongs to the source decomposition; it is distinct from the
morsel grain used to partition an already materialized input.

\evidence{\src{include/ibex/runtime/lazy_table.hpp},
\code{LazySourceReader::decode_units} and \code{decode};
\src{src/runtime/lazy_table.cpp}, \code{acquire_reader} and
\code{release_reader}; \src{src/runtime/pipeline_executor.cpp},
\code{PipelinedScanOperator}.}

\subsection{Map operations: sharing, evaluating, and gathering}
A projection or rename can change a chunk's column entries without visiting
its value rows. The \code{map_chunk} kernel receives source column positions
and output names and constructs entries sharing the selected column values.
The caller supplies the derived row-layout properties. This is a metadata
operation over the columns, rather than a copy of every selected value.

A filter changes the row set and therefore has different work to do. The
predicate identifies survivors; gathering constructs output columns for those
positions, together with their validity. An update evaluates expressions to
produce replacement or additional columns. A fused filter/project step can
gather only the columns required by its output. The physical map description
records kernel capability and column representation/null-policy signatures,
so representation-specific work is chosen at the execution boundary.

Fixed-width kernels can borrow \code{ColumnView} objects containing a value
pointer, a length, and a validity pointer. These views own no storage: the
input chunk or table must remain alive during their use. A string or packed
boolean column requires its own access pattern, rather than being treated
as an ordinary array of independently addressable scalar objects.

\evidence{\src{src/runtime/kernel_types.hpp}, \code{ColumnView},
\code{ChunkView}, and \code{map_chunk};
\src{include/ibex/runtime/pipeline.hpp}, \code{MapStep},
\code{MapKernelCapability}, and \code{ColumnKernelSignature};
\src{tests/test_physical_plan.cpp}, projection and fused-filter tests.}

\subsection{Hash joins: build, probe, and output assembly}
For an eligible equality join, the build phase constructs a lookup structure
over one input. The probe phase looks up keys from the other input and gathers
the matching payload columns. These phases have different lifetime needs:
the index and its backing table must survive every probe, while probe chunks
can be processed one at a time.

\paragraph{Choosing the build side.}
In the ordinary single-key streaming path, the right input is already
materialized. If it is sufficiently small, the engine builds its index without
draining the left input. Otherwise it buffers left chunks until their row count
reaches the right-side count or the left input ends. If the count reaches the
right-side count, it builds on the right and replays the buffered chunks before
continuing the left stream. Only the exhausted smaller-left case needs to
concatenate those chunks to consider building on the left. An ordering benefit
can still favor building on the right. Thus build-side selection depends on
both input size and useful ordering; it is not simply ``always hash the
smaller table''.

\paragraph{Representing repeated keys.}
\code{JoinHashIndex} stores a head row for each key and a \code{chain_next}
array linking build rows with the same key. A uniqueness flag identifies the
case where chain traversal can be avoided. In the serial build, reverse row
insertion makes duplicate chains traverse in original build-row order.
Null keys are excluded in this streaming path, which implements the
\code{nulls never} case. Other join semantics select other implementations.
For categorical probes, the engine can resolve dictionary strings to index
heads once per dictionary and then use row codes to access those heads.

\paragraph{Producing rows.}
Probing produces pairs of left and right row indices. Output assembly gathers
the requested payloads through those pairs and applies the planned output
names. Repeated keys can expand one probe row into many result rows, so memory
for matches depends on result cardinality as well as input cardinality.
Worker-private probe state retains shared ownership of the immutable build.
The orientation and ordering logic must be considered when describing output
order; the order of hash buckets is not a semantic ordering rule.

\paragraph{Filtering a deferred probe source.}
A separate deferred path obtains build keys before decoding the probe payload.
It publishes a dynamic key filter and marks the filter ready. A supported
source can use that filter to narrow candidate rows; exact join probing still
determines matches. The two-phase path can find matching key rows first and
materialize payload columns for survivors afterwards. If that optimization
declines, execution returns to the ordinary materialization/probe path. This
is why a deferred probe source cannot always be pulled independently before
the join's build phase has run.

\evidence{\src{src/runtime/join_chunked.cpp},
\code{ChunkedInnerJoinOperator::initialize}, \code{BufferedThenStreamSource},
\code{choose_and_build_single_key}, \code{JoinHashIndex},
\code{fill_partitioned_heads}, \code{JoinProbe::resolve_categorical_heads},
\code{JoinProbe::assemble_output}, \code{resolve_deferred_probe}, and
\code{try_two_phase_probe}; \src{src/runtime/join_chunked_internal.hpp},
\code{JoinProbeFactory}. This account covers the inspected streaming equality
join family, not every join kind.}

\subsection{Aggregation: state survives the input chunks}
Aggregation replaces many input rows with group results. In the hash path,
group state survives across chunk pulls, while each input chunk is released
after its contribution has been accumulated. The coordinator enforces four
phases:
\begin{enumerate}
  \item \emph{Discovery} identifies the group for each input row and retains
    the active chunk while its column references are in use.
  \item \emph{Accumulation} consumes the discovered group identifiers and
    updates aggregate state, then releases the chunk. Some specialized paths
    combine discovery and accumulation; an explicit transfer marker prevents
    accumulating those values twice.
  \item \emph{Final ordering}, after input exhaustion, brings partition-owned
    groups into global first-occurrence order. Strategies that already use
    global group identifiers need no deferred ordering merge.
  \item \emph{Emission} constructs result columns from finalized group state.
    The implementation rejects emission before input consumption or ordering
    finalization is complete.
\end{enumerate}
For a grouped sum, the conceptual state is an accumulator for each group;
processing row $i$ updates the accumulator selected by its group identifier.
The implementation specializes discovery and accumulation by key and value
representation. Distinct counting needs additional per-group value state,
so ``one accumulator per group'' is not a memory model for every aggregate.
The hash path can consume chunks without retaining all their original rows,
but it emits its result only after the input has ended.

\paragraph{When sorted input changes the algorithm.}
The adaptive aggregate examines the first nonempty chunk and its ordering
metadata. The sorted path requires the leading ordering keys to cover the
grouping columns and rejects grouping entries carrying validity bitmaps.
There are also aggregate-specific gates: distinct counting and nonnumeric
\code{first}/\code{last}, for example, use the hash fallback in this path.
On fallback, the inspected first chunk is prepended to the remaining stream;
it is neither discarded nor read again from the source.

When the sorted strategy is admitted, equal keys occupy contiguous runs.
The operator retains the current group's state across chunk boundaries,
closes it when the key changes, and buffers completed groups for output.
At end of input it closes the remaining open group. The relevant ordering
is ordering by the grouping keys: a time-ascending table with interleaved
symbols does not automatically qualify for aggregation by symbol. This is
an operational reason to distinguish a TimeFrame's time order from a
group-contiguous layout.

\evidence{\src{src/runtime/aggregate_chunked.cpp},
\code{HashAggregateState::next_discovery}, \code{accumulate_discovery},
\code{finalize_ordering}, \code{emit_output},
\code{HashAggregatePhaseOperator::next}, and
\code{ChunkedSortedAggregateOperator}, especially \code{decide_strategy},
\code{sorted_on_group_by}, \code{needs_hash_fallback}, and \code{next_sorted}.}

\subsection{Interchange and memory lifetime}
The Arrow C Data Interface supplies array/schema descriptors and a release
callback for sharing columnar data within a process~\cite{arrowCData}.
Ibex's borrowing import copies values. Its adopting import instead creates a
shared owner for the producer's array and, after successful validation and
import, clears the caller's array descriptor. The retained owner's destructor
eventually calls the producer's release function. Failed adoption leaves
ownership with the caller.

Whether a payload is adopted or converted depends on representation. The
inspected implementation can borrow supported primitive buffers and validated
UTF-8 buffers through that owner. Narrow integer inputs are widened into
Ibex's integer representation; timestamps requiring a unit conversion also
need new value storage. Sharing an external buffer therefore preserves a
producer lifetime, while converting it allocates an Ibex representation.
The implementation's \code{import_column} dispatch is the evidence here:
the public header's shorter adoption comment lists fewer supported layouts.

\evidence{\src{src/interop/arrow_c_data.cpp}, \code{AdoptedArrayOwner},
\code{import_column}, \code{import_string_column},
\code{import_table_from_arrow}, and \code{adopt_table_from_arrow}.}

For operational reasoning, the following lists retained state rather than
asserting an unmeasured whole-query memory bound. Shared buffers must be
counted once even when several entries refer to them.
\begin{center}
\begin{tabularx}{\linewidth}{@{}lX@{}}
\toprule
Boundary & Storage that may remain live\\
\midrule
Lazy binding & Whole-source columns already decoded into its cache.\\
Map pipeline & Input buffers, worker scratch/output, and buffered chunks.\\
Hash join & Build table and index; buffered probe input and match indices.\\
Hash aggregate & Group keys and aggregate state across all input chunks.\\
Sorted aggregate & Open-group state and buffered completed-group output.\\
Result sink & The growing result and the current incoming chunk.\\
Arrow adoption & Producer storage retained by the shared array owner.\\
\bottomrule
\end{tabularx}
\end{center}
The distinction between input consumption and output production is central:
a bounded chunk interface can feed an operator whose retained state grows
with the number of groups, join matches, or result rows.

\section{A trade-summary workload}
\label{sec:trade-workload}
The runnable example \src{docs/engine/trade_summary.ibex} starts with five
timestamped ticks in ascending time order. It constructs a TimeFrame, removes
a zero-price tick, computes trade notionals, enriches symbols with an instrument
table, and produces volume-weighted average prices (VWAP). This first workload
uses a complete sample interval; window and resampling rules remain deferred.

\lstinputlisting[firstline=15]{trade_summary.ibex}

The instrument table has one row per symbol, with venue \code{XNAS} for both
symbols. The expected data flow is:
\begin{center}
\begin{tabularx}{\linewidth}{@{}lrlX@{}}
\toprule
Stage & Rows & Operation & Resulting work\\
\midrule
Ticks & 5 & Construct & Timestamp, symbol, price, volume columns.\\
Eligible & 4 & Filter/update & Drop the zero price; calculate price times volume.\\
Enriched & 4 & Join & Add the matching venue; these keys do not multiply rows.\\
Totals & 2 & Aggregate & Accumulate volume and notional by symbol and venue.\\
Report & 2 & Update/order & Divide totals and explicitly sort by symbol.\\
\bottomrule
\end{tabularx}
\end{center}
The report contains AAPL with volume $30$, notional $3040$, and VWAP
$3040/30$; MSFT has volume $20$, notional $4030$, and VWAP $201.5$.
Keeping the division after aggregation expresses
$\sum_i p_i v_i / \sum_i v_i$, rather than an unweighted mean of prices.
The explicit final order makes this summary's presentation independent of
the first occurrence of each symbol in the ticks.

\paragraph{Observed execution boundary.}
Running the example with the existing release executable and
\code{--no-history} and \code{--report-planner} produced those two result rows and
reported statement execution because there were no lazy table sources to plan
against. Named intermediate results in this run are therefore separate
statement results. The sequence above describes data dependencies; it is
not a claim that the entire script became one fused physical pipeline.
The numeric output validates this small workload, not the parallel strategy
gates or compiled/interpreted parity. No engine rebuild or benchmark was
performed for this check.

Replacing the constructed input with a lazy file source would add demand
analysis and decoder behavior to the trace. Its exact physical plan should
be recorded from that run before making claims about pushdown or fusion.
Similarly, adding windows later will introduce persistent ordered state whose
boundaries need a separate treatment.

\section{Execution cost and existing measurements}
\label{sec:cost}
Ibex already has a substantial measurement corpus. The scale suite and its
published benchmark page provide observations across operators, row counts,
and thread configurations~\cite{ibexBenchmarks}. This section uses a small,
explicitly identified snapshot to ground the cost explanation. It separates
what those measurements show from the mechanisms that can explain a cost;
timing a query alone does not identify its bottleneck.

\subsection{What determines the amount of work?}
Row count is only one dimension. Let $n$ denote input rows, $s$ the fraction
surviving a filter, $w$ the bytes of demanded payload per row, $g$ the number
of groups, and $m$ the number of join-result rows. These are explanatory
quantities, not parameters of a claimed Ibex optimizer cost model.

\paragraph{Reading and selecting.}
Column demand limits which payloads must be read. A selective scan can reduce
payload materialization toward $snw$ bytes, but predicate evaluation still
has its own input cost. Source encodings and page boundaries determine whether
skipped rows also avoid decoding work. A small result therefore does not imply
a proportionally cheap scan. Cached columns change the work again: gathering
from existing buffers avoids source decoding but still reads and writes values.

\paragraph{Computing and moving values.}
A row-local expression such as doubling a price visits $n$ values and produces
$n$ outputs. A projection can share payloads, whereas a filter generally gathers
survivors into output buffers. Fusion can remove intermediate passes and avoid
gathering unused columns. Fixed-width output volume depends on rows and width;
string output also depends on the number of character bytes. Consequently,
operations with similar arithmetic can have different allocation and copying
costs. These mechanisms motivate a bandwidth hypothesis, but the published
wall times alone do not establish memory-bandwidth saturation.

\paragraph{Maintaining state.}
Aggregation combines row processing with state indexed by $g$ groups. More
groups, additional key components, and different key representations alter
lookup and state-management work. A hash join similarly pays for its build
index, probe lookups, and materializing $m$ matches. Duplicate keys can make
$m$ much larger than either input. Reducing output width helps the gather,
but does not eliminate the cost of identifying the matches.

\paragraph{Ordering and reducing output.}
Sorting a complete table and selecting its top 100 rows request different
amounts of output and can use different algorithms. Ibex's bounded-heap top-k
path performs selection before gathering its small result, rather than
materializing a full sorted table. Useful input ordering can also avoid
sorting or admit the sorted aggregate described in Section~\ref{sec:operations}.
Maintaining time order can thus have an execution benefit as well as a semantic
purpose, although time order and group-contiguous order are different facts.

\evidence{\src{src/runtime/lazy_table.cpp}, selected decoding;
\src{include/ibex/runtime/pipeline.hpp}, fused map descriptions;
\src{src/runtime/join_chunked.cpp} and
\src{src/runtime/aggregate_chunked.cpp}, state and output construction;
\src{src/runtime/runtime_entry.cpp}, the TopK execution-path description.
These are implementation-based explanations, not measured attribution of the
timings below to individual phases.}

\subsection{The measurement boundary}
The inspected \src{docs/benchmarks.html} snapshot was generated on
3 September 2026 at 13:08 UTC and records commit
\code{0d991d3677f2927eb40c77ba9f489ad6877baa24}. Its source is
\src{benchmarking/results/scales_aws_20260903T100845.csv}.
The published methodology describes AWS r7i.2xlarge instances with 8 vCPUs
(4 physical cores with two hardware threads each), 64 GB of memory, and one
instance per engine. The timings below are reported means from that run,
not fresh measurements of the current checkout.

For that archived snapshot, the Ibex wrapper selects resident input tables
and parses and lowers the query before the timed iterations. It warms up, then times
\code{runtime::interpret}, including construction of the materialized result.
This measures execution rather than file loading or C++ compilation. Separate
parse-inclusive and scan-inclusive harness modes have different boundaries
and must not be mixed into that interpretation. Physical execution decisions
made inside \code{interpret} remain inside the timer.
The revised in-memory harness includes parsing and lowering; its new results
must be distinguished from this historical snapshot.

\evidence{\src{docs/benchmarks.html}, embedded \code{PAYLOAD} metadata and
timings; \src{docs/methodology.html};
\src{benchmarking/gen_website.py}, \code{avg_ms} payload construction;
\src{tools/ibex_bench.cpp}, \code{run_benchmark}, especially the
\code{!include_parse} branch. The selected page cells were checked against
their source CSV. The page snapshot does not expose individual iteration
samples or the invocation's repetition count; harness defaults alone do not
establish that count for this archived run.}

\subsection{Representative observations}
All times in Table~\ref{tab:cost-snapshot} are milliseconds. ``Default'' is
the published \code{ibex} configuration; ``one'' is \code{ibex-st}, the same
build with \code{IBEX_CORES=1}. Query identifiers are retained so the exact
expressions can be found in the methodology page and harness.

\begin{table}[htbp]
\centering
\small
\begin{tabular}{@{}lrrr@{}}
\toprule
Query & 1M, default & 16M, default & 16M, one\\
\midrule
\code{update_price_x2} & 0.212 & 4.363 & 18.995\\
\code{mean_by_symbol} & 0.729 & 7.424 & 30.339\\
\code{sort_price} & 21.431 & 516.333 & 697.558\\
\code{order_head_topk} & 2.650 & 39.369 & 39.481\\
\code{tf_rolling_sum_1m} & 2.880 & 53.370 & 50.783\\
\code{tf_rolling_median_1m} & 29.884 & 485.317 & 483.917\\
\code{tf_resample_1m_ohlc} & 4.342 & 98.409 & 96.602\\
\bottomrule
\end{tabular}
\caption{Selected means from the identified September 2026 benchmark
snapshot~\cite{ibexBenchmarks}. M denotes one million input rows.}
\label{tab:cost-snapshot}
\end{table}

Three observations are useful for understanding the engine. First, at 16M
rows, \code{mean_by_symbol} emits 252 groups, \code{order_head_topk} emits
100 rows, and \code{sort_price} emits all 16M rows. Their timings concern
different work and result sizes. In particular, top-k and full sorting are
not interchangeable benchmark answers: the harness requests the top 100
descending prices for the former and a full ascending sort for the latter.

Second, comparing the reported 16M means gives about $4.35\times$ acceleration
for doubling prices and $4.09\times$ for mean by symbol, but only about
$1.35\times$ for full price sorting and essentially none for this top-k
query. These ratios describe the complete operations at that scale. They do
not mean that every phase has that speedup, or that eight vCPUs are eight
independent physical cores.

Third, the time-series rows expose substantial differences within the same
workload family: the one-minute rolling median takes 485.317 ms at 16M rows,
versus 53.370 ms for the rolling sum. The default and one-core means are close
for these examples and for resampling. That is evidence of little observed
threading benefit in this snapshot, not a statistical verdict about their
small differences. The harness constructs timestamps at one-second spacing
and the expressions include \code{as_timeframe}; these numbers therefore do
not isolate window-state maintenance. Varying timestamp density or window
duration would change a different dimension from simply adding more rows.
The detailed window algorithms remain a later investigation.

\evidence{The source CSV's \code{rows} and \code{avg_ms} columns;
\src{tools/ibex_bench.cpp}, definitions of \code{order_head_topk},
\code{sort_price}, and the TimeFrame benchmark group. The published page also
notes that its cross-engine rolling-EWMA comparison uses different mathematics;
that row is not treated here as a comparison of equivalent computations.}

\subsection{Wall time, CPU work, and memory answer different questions}
Parallel execution overlaps work, so adding per-operator or per-worker times
does not generally reconstruct elapsed query time. The critical dependencies
include build-before-probe, aggregate finalization, and ordered output delivery.
Task submission, buffering, and merging also cost time. A larger thread budget
is useful only where it reduces elapsed time enough to offset those costs.

The repository's \src{benchmarking/work_multiplication.py} complements wall
time with total child-process CPU consumption at different core counts.
\src{MEASURING.md} records why this matters: a historical sharded decoder
repeated work while still showing attractive parallel occupancy. High CPU
growth is a diagnostic signal, not proof of redundant decoding; contention
can also increase CPU time. Counting rows, calls, or bytes is needed to
distinguish those explanations.

The benchmark's memory column is absolute peak resident memory, including
already-resident input and other live process storage. It is not an
operator's incremental allocation count. The Ibex harness attempts to reset
Linux's high-water mark before timed iterations; if that mechanism is
unavailable, the code documents a lifetime-peak fallback. Consequently, these
figures describe the footprint of the recorded harness execution and cannot
be substituted directly for the per-operator state accounting above.

\evidence{\src{MEASURING.md}, work multiplication and profile interpretation;
\src{benchmarking/work_multiplication.py}; \src{tools/ibex_bench.cpp},
\code{reset_peak_rss} and \code{peak_rss_mb}.}

\subsection{Using the wider corpus}
The scale suite provides broad coverage, while the PDS-H query suite,
operator profiles, and focused time-series experiments provide complementary
evidence. Their measurement boundaries and datasets should be recorded before
combining results. Existing measurements are the starting point; new runs are
needed only for a specific unresolved claim, such as the contribution of a
particular decoder or the cost of maintaining an ordering guarantee.

For causal claims about a change, the repository's paired, interleaved A/B
workflow and per-operator profiles are more appropriate than differences
between independently published snapshots. This document has reused the
archived evidence without running a benchmark. It neither extrapolates those
times to other machines nor treats a single snapshot as a current performance
guarantee.


\section{Agreed direction and next investigations}
Author review established the following direction for developing this
foundation. The entries distinguish agreed priorities from investigations
that have yet to be completed.
\begin{enumerate}
  \item \textbf{Intended emphasis: agreed.} Address an external
    database-systems audience, developing design rationale and literature
    comparisons alongside the operational detail needed to understand
    execution.
  \item \textbf{A representative workload: time series.} Develop a verified
    time-series example as the document's running workload. The author's
    motivation is that this area deserves more attention in analytics-engine
    design; this is a choice of emphasis, not a claim of demonstrated
    superiority over other engines. Section~\ref{sec:trade-workload} establishes
    a small trade-summary workload to extend in later revisions.
  \item \textbf{Design intent: ordering and numerical behavior.} Preserve
    output row-order guarantees because they are valuable for time-series
    reasoning. Floating-point reproducibility requirements may be weakened
    independently, with their precise scope deferred to later work.
    Functional equivalence of the interpreted and compiled routes is an
    established requirement. Any intentional exception should be stated
    explicitly and distinguished from the parity bugs listed below. Whether
    the single-query lease is a long-term constraint remains open.
  \item \textbf{Operator algorithms: initial account added.}
    Section~\ref{sec:operations} traces streaming equality joins, aggregate
    phases, and sorted-input gates. Broader join kinds, key-layout
    specializations, and their algorithm-specific literature still need audit.
  \item \textbf{Time-series execution: deferred.} Later, trace window boundaries, ordered state,
    group layout, resampling, and null handling against both the specification
    and executable examples. Choosing a time-series running workload does
    not require completing this detailed chapter immediately.
  \item \textbf{Data boundaries and memory: initial account added.}
    Section~\ref{sec:operations} traces lazy projection, selected decoding,
    reader ownership, Arrow adoption, and retained state. Source-specific
    decoder behavior, cross-chunk dictionary compatibility, and measured
    peak memory remain open; a shared interface is insufficient evidence for
    every plugin boundary.
  \item \textbf{Cost explanation and existing measurements: added.}
    Section~\ref{sec:cost} explains cost drivers using an identified published
    benchmark snapshot. Precise numerical reproducibility rules and new
    performance experiments remain deferred. Further measurement should target
    gaps in the existing corpus and follow \src{MEASURING.md}.
\end{enumerate}

\appendix
\section{Interpreted/compiled parity bug list}
\label{sec:backend-bugs}
This list records functional divergences encountered while preparing the
document. Entries are open implementation bugs under the equivalence contract,
unless an explicit design decision identifies an intentional exception.
No such exception has been established for the entries below. The list is
not an exhaustive parity audit, and these identifiers are local to this
document rather than external issue numbers.

\begin{description}
  \item[PARITY-001: Window selection.] The Window case in
    \code{Emitter::emit_node} throws when \code{win.select_only()} is true.
    Expected behavior: compile \code{window + select} with the same semantics
    as interpreted execution, including the selected output schema.
  \item[PARITY-002: Aligned windows.] The same case throws when
    \code{win.aligned()} is true. Expected behavior: compiled execution
    implements the same epoch-aligned window boundaries and results.
  \item[PARITY-003: Windowed tuple updates.] The same case rejects nonempty
    \code{upd.tuple_fields()}. Expected behavior: windowed tuple assignments
    have equivalent functionality in both routes. The existing test
    ``emitter: window node - tuple fields are still rejected'' records the
    current rejection; it does not make that rejection an intentional
    semantic exception.
  \item[PARITY-004: Model clauses.] The Model case unconditionally throws
    ``model clause is not yet supported in compiled mode''. Expected behavior:
    compiled execution supports the model-clause functionality provided by
    interpreted execution.
\end{description}

\evidence{\src{src/codegen/emitter.cpp}, \code{Emitter::emit_node}, Window
and Model cases; \src{tests/test_codegen.cpp}, the window tuple-rejection
test. \src{SPEC.md} defines window selection, aligned windows, and tuple
assignments. These findings are based on source inspection; end-to-end
reproductions and fixes remain outstanding.}

For each newly encountered divergence, record the triggering program or IR
shape, expected shared behavior, observed difference, and evidence. Distinguish
source-inspected findings from reproduced failures. Close an entry after a fix
and a regression check comparing the two routes, or record the explicit design
decision that makes it an intentional exception.

\section{Source map for continued investigation}
Paths are relative to the repository root. Symbols and test names above are
preferred to line numbers because this checkout is evolving.
\begin{description}
  \item[Language to IR] \src{src/parser/lower.cpp};
    \src{src/ir/optimizer.cpp}; \src{src/ir/canonicalize.cpp}.
  \item[Caller-specific planning] \src{src/repl/repl.cpp};
    \src{src/ir/required_columns.cpp}; \src{src/ir/scan_predicates.cpp}.
  \item[Execution entry and decisions] \src{src/runtime/interpreter.cpp};
    \src{src/runtime/runtime_entry.cpp}; \src{src/runtime/physical_plan.hpp};
    \src{src/runtime/physical_plan.cpp}.
  \item[Flow and scheduling] \src{include/ibex/runtime/operator.hpp};
    \src{src/runtime/pipeline_executor.cpp}; \src{src/runtime/worker_pool.cpp}.
  \item[Representation and metadata] \src{include/ibex/core/column.hpp};
    \src{include/ibex/runtime/interpreter.hpp};
    \src{include/ibex/runtime/table_properties.hpp}.
  \item[Next operator investigations] \src{src/runtime/join_chunked.cpp};
    \src{src/runtime/aggregate_chunked.cpp}; \src{src/runtime/window.cpp};
    \src{src/runtime/ordered_chunked.cpp}. These are investigation targets,
    not evidence of a complete audit of every strategy in these files.
  \item[Backend and regression evidence] \src{src/codegen/emitter.cpp};
    \src{tests/test_physical_plan.cpp}; \src{tests/test_table_properties.cpp};
    \src{tests/test_column.cpp}.
\end{description}

\bibliographystyle{plain}
\bibliography{references}
\end{document}
